
\documentclass{article}
\usepackage{amsmath,amssymb,graphicx}
\usepackage[margin=1in]{geometry}

\title{Summary of ``An Extreme Learning Machine-Based Method for Computational PDEs in Higher Dimensions''}
\author{Prepared by ChatGPT}
\date{\today}

\begin{document}
\maketitle

\section{Overview}
The paper by Y.~Wang and S.~Dong (Computer Methods in Applied Mechanics and Engineering, 418:116578, 2024) proposes two \emph{randomised neural network} solvers for high–dimensional partial differential equations (PDEs):
\begin{enumerate}
  \item \textbf{Extreme Learning Machine (ELM)} with fixed hidden--layer parameters;
  \item \textbf{ELM/A--TFC}, which embeds an ``Approximate'' Theory of Functional Connections (A--TFC) constraint to enforce boundary/initial conditions exactly while leaving a free field represented by an ELM network.
\end{enumerate}
Both techniques transform the PDE residual at randomly sampled collocation points into a (non‑)linear least–squares system for the \emph{output} weights, avoiding gradient‐based training on deep architectures.

\section{Algorithmic Details}

\subsection{Notation}
Let $\Omega\subset\mathbb{R}^d$ be a rectangular domain and consider the generic boundary‑value (or initial–boundary‑value) problem
\begin{align}
\mathcal{L}u(x)+\mu \mathcal{N}(u(x)) &= Q(x), && x\in\Omega,\\
\mathcal{B}u(x) &= H(x), && x\in\partial\Omega,
\end{align}
with possibly time‑dependent operators.  Here $\mathcal{L}$ is linear, $\mathcal{N}$ may be nonlinear, and $\mathcal{B}$ represents boundary/initial conditions.

\subsection{ELM solver}
\begin{enumerate}
\item \textbf{Random projection.}  Choose the activation $g(\cdot)$ (e.g.\ $\tanh$ or $\sin$).  Draw weights $w_j\in[-R_m,R_m]^d$ and biases $b_j\in[-R_m,R_m]$ independently and fix them forever.
\item \textbf{Ansatz.}  Represent the solution by a single‑hidden‑layer network
\(
u_M(x)=\sum_{j=1}^{M} a_j\,g(w_j\!\cdot\! x + b_j),
\)
where only the output weights $a_j$ are trainable.
\item \textbf{Collocation sampling.}  Generate
  \begin{itemize}
    \item $N_{\text{in}}$ interior points $\{x^{(k)}\}$,
    \item $N_{\text{bc}}$ boundary points $\{x^{(k)}_{\partial}\}$ (and $N_{t_0}$ initial points for parabolic PDEs).
  \end{itemize}
\item \textbf{Residual assembly.}
  \begin{align*}
  &\text{Interior: } \quad r^{(k)} = \mathcal{L}u_M(x^{(k)})+\mu\mathcal{N}(u_M(x^{(k)}))-Q\bigl(x^{(k)}\bigr),\\
  &\text{Boundary: } \quad b^{(k)} = \mathcal{B}u_M(x^{(k)}_{\partial})-H\bigl(x^{(k)}_{\partial}\bigr).
  \end{align*}
\item \textbf{Least‑squares training.}  Form the stacked residual vector $F(a)\in\mathbb{R}^{N}$ ($N=N_{\text{in}}+N_{\text{bc}}$).\\[2pt]
  \begin{itemize}
    \item \emph{Linear PDE}: $F(a)=Aa-q$ is linear in $a$; solve $a^\star=\arg\min\|Aa-q\|_2$ via pseudo‑inverse.
    \item \emph{Nonlinear PDE}: apply Levenberg–Marquardt to $\min_a \|F(a)\|_2^2$ (typically converges within $\mathcal{O}(10)$ iterations owing to small parameter count).
  \end{itemize}
\item \textbf{Prediction.}  The trained network $u_M(x;a^\star)$ approximates $u(x)$ everywhere.
\end{enumerate}

\subsection{ELM/A--TFC solver}
\begin{enumerate}
  \item Construct an \emph{A--TFC constrained expression}
  \(
  u(x)=G(x)+\mathcal{C}[f](x),
  \)
  where $G$ is a simple function satisfying the non‑homogeneous boundary data and $\mathcal{C}[\,\cdot\,]$ embeds homogeneous boundary conditions.
  \item Approximate the free function $f(x)$ by an ELM network with coefficients $\alpha_j$.
  \item Follow steps 3--5 above to train $\alpha_j$.  By construction the final $u(x)$ exactly satisfies the boundary/initial conditions.
\end{enumerate}

\section{Theoretical Findings}

Building on the work of Rahimi \& Barron and, specifically, the result by Igelnik \& Pao (1995), the authors recall that for any Lipschitz $f$ on $[0,1]^d$ there exists a random‑feature expansion with \emph{expected} $L^2$ error
\(
\mathbb{E}\,\|f-f_M\|_{L^2}\le C\,M^{-1/2},
\)
where $C$ depends on the Lipschitz constant but \emph{not} on the dimension $d$.  This dimension‑independent $M^{-1/2}$ rate motivates using large but \emph{shallow} random networks for high‑dimensional PDEs.

\section{Numerical Results}

The paper reports extensive experiments for dimensions $d=3$--$10$ using networks of width $M=O(10^3)$.  Below we list the headline results:

\begin{itemize}
  \item \textbf{Poisson equation} on $[-1,1]^d$ ($d=3,5,7,9$): ELM achieved $L^\infty$ errors from $2.3\times10^{-10}$ (3‑D) to $4.0\times10^{-5}$ (9‑D) with training times $2.6$--$14.6$\,s---\emph{two orders of magnitude faster and $10^2$--$10^4$ times more accurate} than a comparable PINN (Table~16).
  \item \textbf{Non‑linear Poisson}: similar trends; ELM errors $\le1.6\times10^{-4}$ (9‑D) vs.\ PINN $\approx 5\times10^{-2}$, with $\sim10\times$ shorter training (Table~17).
  \item \textbf{Advection–diffusion ($d=3,6,10$)}: exponential error decay w.r.t.\ hidden width until saturation; max error $<10^{-7}$ for $d=3$ (Fig.~11).  Optimal sampling radius $R_m$ decreases with $d$ (Tables~5–6).
  \item \textbf{Heat equation ($d=7$)} solved by \emph{ELM/A--TFC}: point‑wise errors $\mathcal{O}(10^{-4})$ and exponential convergence versus boundary points (Figs.~24–25).
\end{itemize}

Overall, the ELM family delivers PINN‑level (or better) accuracy at drastically lower computational cost, and the dimensional curse is mitigated but not entirely removed.

\section{Conclusion}
Randomised, shallow ELM networks—trained by one‐shot least‑squares—constitute a simple yet powerful alternative to gradient‑based deep PINNs for high‑dimensional PDEs.  The optional A--TFC embedding further improves boundary fidelity at modest extra cost.

\end{document}
